\subsection{Null Spaces and Ranges}

\begin{exercise}[3b1]
    Give an example of a linear map $T$ with $\dim \nul T = 3$ and $\dim \range T = 2$.
\end{exercise}

\begin{sol}
    Let $T \in \lin{\r^5, \r^2}$ be defined by
    \[T(x_1, x_2, x_3, x_4, x_5) = (x_1, x_2).\]
    Then $\nul T = \set{(0, 0, x_3, x_4, x_5) \mid x_3, x_4, x_5 \in \r}$, which has dimension $3$, and $\range T = \r^2$, which has dimension $2$.
\end{sol}

\begin{exercise}[3b2]
    Suppose $S, T \in \lin{V}$ are such that $\range S \subseteq \nul T$.
	Prove that $(ST)^2 = 0$.
\end{exercise}

\begin{sol}
    Let $v \in V$.
	Then
    \[(ST)^2 v = ST(STv) = S(T(STv)).\]
    Since $STv \in \range S$ and $\range S \subseteq \nul T$, we have $T(STv) = 0$.
	Therefore,
    \[(ST)^2 v = S(0) = 0.\]
    Since $v$ was arbitrary, it follows that $(ST)^2 = 0$.
\end{sol}

\begin{exercise}[3b3]
    Suppose $v_1, \ldots, v_m$ is a list of vectors in $V$.
	Define $T \in \lin{\f^m, V}$ by
    \[T(z_1, \ldots, z_m) = z_1 v_1 + \cdots + z_m v_m.\]
    \begin{subexercises}
        \item What property of $T$ corresponds to $v_1, \ldots, v_m$ spanning $V$?
        \item What property of $T$ corresponds to the list $v_1, \ldots, v_m$ being linearly independent?
    \end{subexercises}
\end{exercise}

\begin{sole}
    \item The list $v_1, \ldots, v_m$ spans $V$ if and only if $\range T = V$, i.e., $T$ is surjective.
    \item The list $v_1, \ldots, v_m$ is linearly independent if and only if $\nul T = \set{0}$, i.e., $T$ is injective.
\end{sole}

\begin{exercise}[3b4]
    Show that $\set{T \in \lin{\r^5, \r^4} \mid \dim \nul T > 2}$ is not a subspace of $\lin{\r^5, \r^4}$.
\end{exercise}

\begin{sol}
    Let $T_1, T_2 \in \lin{\r^5, \r^4}$ be defined by
    \[T_1(x_1, x_2, x_3, x_4, x_5) = (x_1, x_2, 0, 0)\]
    and
    \[T_2(x_1, x_2, x_3, x_4, x_5) = (0, 0, x_3, x_4).\]
    Then $\nul T_1 = \set{(0, 0, x_3, x_4, x_5) \mid x_3, x_4, x_5 \in \r}$ and $\nul T_2 = \set{(x_1, x_2, 0, 0, x_5) \mid x_1, x_2, x_5 \in \r}$ both have dimension $3$, so $T_1$ and $T_2$ are in the set $\set{T \in \lin{\r^5, \r^4} \mid \dim \nul T > 2}$.
	However,
    \[(T_1 + T_2)(x_1, x_2, x_3, x_4, x_5) = (x_1, x_2, x_3, x_4),\]
    which has null space $\nul(T_1 + T_2) = \set{(0, 0, 0, 0, x_5) \mid x_5 \in \r}$ of dimension $1$.
	Therefore,
    \[\dim \nul(T_1 + T_2) = 1 < 2,\]
    so $T_1 + T_2$ is not in the set.
	Hence, the set is not closed under addition and is not a subspace of $\lin{\r^5, \r^4}$.
\end{sol}

\begin{exercise}[3b5]
    Give an example of $T \in \lin{\r^4}$ such that $\range T = \nul T$.
\end{exercise}

\begin{sol}
    Let $v_1, v_2, v_3, v_4$ be the standard basis of $\r^4$.
	Define $T \in \lin{\r^4}$ by
    \[Tv_1 = v_3, \quad Tv_2 = v_4, \quad Tv_3 = 0, \quad Tv_4 = 0.\]
    Then for any $x = (x_1, x_2, x_3, x_4) \in \r^4$,
    \[Tx = x_1 Tv_1 + x_2 Tv_2 + x_3 Tv_3 + x_4 Tv_4 = x_1 v_3 + x_2 v_4.\]
    Then $\range T = \spn{v_3, v_4}$.
	If $Tx = 0$, then $x_1 v_3 + x_2 v_4 = 0$, which implies $x_1 = x_2 = 0$.
	Therefore, $\nul T = \set{(0, 0, x_3, x_4) \mid x_3, x_4 \in \r} = \spn{v_3, v_4}$.
	Hence, $\range T = \nul T$.
\end{sol}

\begin{exercise}[3b6]
    Prove that there does not exist $T \in \lin{\r^5}$ such that $\range T = \nul T$.
\end{exercise}

\begin{sol}
    Suppose for contradiction that there exists $T \in \lin{\r^5}$ such that $\range T = \nul T$.
	Then by the fundamental theorem of linear maps,
    \[\dim \r^5 = \dim \nul T + \dim \range T = 2\dim \range T.\]
    Therefore, $5 = 2\dim \range T$, which is impossible since $5$ is odd and $2\dim \range T$ is even.
	Hence, no such $T$ exists.
\end{sol}

\begin{exercise}[3b7]
    Suppose $V$ and $W$ are finite-dimensional with $2 \leq \dim V \leq \dim W$.
	Show that $\set{T \in \lin{V, W} \mid T \text{ is not injective}}$ is not a subspace of $\lin{V, W}$.
\end{exercise}

\begin{sol}
    Let $\dim V = n$ and $\dim W = m$ with $2 \leq n \leq m$.
	Let $v_1, v_2, \ldots, v_n$ be a basis of $V$, and let $w_1, w_2, \ldots, w_m$ be a basis of $W$.
	Define $T_1, T_2 \in \lin{V, W}$ by
    \[T_1 v_1 = w_1, \quad T_1 v_k = 0 \text{ for } k \geq 2,\]
    and
    \[T_2 v_1 = 0, \quad T_2 v_k = w_k \text{ for } k \geq 2.\]
    Then $v_2 \in \nul T_1$ and $v_1 \in \nul T_2$, so both $T_1$ and $T_2$ are not injective.
	However,
    \[(T_1 + T_2)v_1 = T_1 v_1 + T_2 v_1 = w_1 + 0 = w_1 \neq 0,\]
    and for $k \geq 2$,
    \[(T_1 + T_2)v_k = T_1 v_k + T_2 v_k = 0 + w_k = w_k \neq 0.\]
    Suppose $v \in V$ is such that $(T_1 + T_2)v = 0$.
	Then $v$ can be expressed as a linear combination of the basis vectors:
    \[v = \alpha_1 v_1 + \alpha_2 v_2 + \cdots + \alpha_n v_n.\]
    Then
    \[(T_1 + T_2)v = \alpha_1 w_1 + \alpha_2 w_2 + \cdots + \alpha_n w_n = 0.\]
    Since $w_1, w_2, \ldots, w_m$ are linearly independent, it follows that $\alpha_1 = \alpha_2 = \cdots = \alpha_n = 0$, hence $v = 0$.
	Therefore, $T_1 + T_2$ is injective.
	Since $T_1$ and $T_2$ are not injective, but their sum is injective, the set $\set{T \in \lin{V, W} \mid T \text{ is not injective}}$ is not closed under addition and thus is not a subspace of $\lin{V, W}$.
\end{sol}

\begin{exercise}[3b8]
    Suppose $V$ and $W$ are finite-dimensional with $\dim V \geq \dim W \geq 2$.
	Show that $\set{T \in \lin{V, W} \mid T \text{ is not surjective}}$ is not a subspace of $\lin{V, W}$.
\end{exercise}

\begin{sol}
    We proceed similarly to the previous exercise.
	Let $\dim V = n$ and $\dim W = m$ with $n \geq m \geq 2$.
	Let $v_1, v_2, \ldots, v_n$ be a basis of $V$, and let $w_1, w_2, \ldots, w_m$ be a basis of $W$.
	Define $T_1, T_2 \in \lin{V, W}$ by
    \[T_1 v_1 = w_1, \quad T_1 v_k = 0 \text{ for } k \geq 2,\]
    and
    \[T_2 v_1 = 0, \quad T_2 v_k = w_k \text{ for } k = 2, \ldots, m, \quad T_2 v_k = 0 \text{ for } k > m.\]
    Then $\range T_1 = \spn{w_1}$ and $\range T_2 = \spn{w_2, \ldots, w_m}$, both of which are proper subspaces of $W$, so $T_1$ and $T_2$ are not surjective.
	However,
    \[(T_1 + T_2)v_1 = T_1 v_1 + T_2 v_1 = w_1 + 0 = w_1,\]
    and for $k = 2, \ldots, m$,
    \[(T_1 + T_2)v_k = T_1 v_k + T_2 v_k = 0 + w_k = w_k.\]
    Therefore, $\range(T_1 + T_2) = \spn{w_1, w_2, \ldots, w_m} = W$, so $T_1 + T_2$ is surjective.
	Since $T_1$ and $T_2$ are not surjective, but their sum is surjective, the set $\set{T \in \lin{V, W} \mid T \text{ is not surjective}}$ is not closed under addition and thus is not a subspace of $\lin{V, W}$.
\end{sol}

\begin{exercise}[3b9]
    Suppose $T \in \lin{V, W}$ is injective and $v_1, \ldots, v_n$ is linearly independent in $V$.
	Prove that $Tv_1, \ldots, Tv_n$ is linearly independent in $W$.
\end{exercise}

\begin{sol}
    Suppose $a_1, \ldots, a_n \in \f$ are such that
    \[a_1 Tv_1 + \cdots + a_n Tv_n = 0.\]
    By linearity of $T$, we have
    \[T(a_1 v_1 + \cdots + a_n v_n) = 0.\]
    Injectivity of $T$ implies that $\nul T = \set{0}$, so
    \[a_1 v_1 + \cdots + a_n v_n = 0.\]
    Since $v_1, \ldots, v_n$ is linearly independent, it follows that $a_1 = \cdots = a_n = 0$.
	Therefore, $Tv_1, \ldots, Tv_n$ is linearly independent in $W$.
\end{sol}

\begin{exercise}[3b10]
    Suppose $v_1, \ldots, v_n$ spans $V$ and $T \in \lin{V, W}$.
	Show that $Tv_1, \ldots, Tv_n$ spans $\range T$.
\end{exercise}

\begin{sol}
    Let $w \in \range T$.
	Then there exists $v \in V$ such that $Tv = w$.
	Since $v_1, \ldots, v_n$ spans $V$, we can write
    \[v = a_1 v_1 + \cdots + a_n v_n\]
    for some scalars $a_1, \ldots, a_n \in \f$.
	By linearity of $T$, we have
    \[w = Tv = T(a_1 v_1 + \cdots + a_n v_n) = a_1 Tv_1 + \cdots + a_n Tv_n.\]
    Therefore, $w$ is a linear combination of $Tv_1, \ldots, Tv_n$.
	Since $w$ was arbitrary in $\range T$, it follows that $Tv_1, \ldots, Tv_n$ spans $\range T$.
\end{sol}

\begin{exercise}[3b11]
    Suppose that $V$ is finite-dimensional and that $T \in \lin{V, W}$.
	Prove that there exists a subspace $U$ of $V$ such that
    \[U \cap \nul T = \set{0} \quad \text{and} \quad \range T = \set{Tu \mid u \in U}.\]
    \remark{Note that this exercise proves the existence of a subspace $U$ of $V$ such that $V = U \op \nul T$.}
\end{exercise}

\begin{sol}
    Since $\nul T$ is a subspace of $V$, we can extend a basis of $\nul T$ to a basis of $V$.
	Let $n_1, \ldots, n_k$ be a basis of $\nul T$, and extend this to a basis of $V$ by adding vectors $u_1, \ldots, u_m$.
	Hence, a basis for $V$ is given by
    \[n_1, \ldots, n_k, u_1, \ldots, u_m.\]
    Define $U = \spn(u_1, \ldots, u_m)$.
	Suppose $v \in U \cap \nul T$.
	Then $v$ can be expressed as a linear combination of the basis vectors of $U$:
    \[v = \alpha_1 u_1 + \cdots + \alpha_m u_m.\]
    Since $v \in \nul T$ as well, we can also express $v$ as a linear combination of the basis vectors of $\nul T$:
    \[v = \beta_1 n_1 + \cdots + \beta_k n_k.\]
    Equating the two expressions for $v$, we have
    \[\alpha_1 u_1 + \cdots + \alpha_m u_m = \beta_1 n_1 + \cdots + \beta_k n_k.\]
    Since $n_1, \ldots, n_k, u_1, \ldots, u_m$ is a basis for $V$, the only way this equality can hold is if all coefficients are zero, i.e., $\alpha_1 = \cdots = \alpha_m = 0$ and $\beta_1 = \cdots = \beta_k = 0$.
	Therefore, $v = 0$, and we conclude that $U \cap \nul T = \set{0}$.

    Now, let $w \in \range T$.
	Then there exists $v \in V$ such that $Tv = w$.
	We can express $v$ as a linear combination of the basis vectors of $V$:
    \[v = \gamma_1 n_1 + \cdots + \gamma_k n_k + \delta_1 u_1 + \cdots + \delta_m u_m.\]
    Applying $T$ to both sides, we have
    \[Tv = T(\gamma_1 n_1 + \cdots + \gamma_k n_k + \delta_1 u_1 + \cdots + \delta_m u_m) = \gamma_1 Tn_1 + \cdots + \gamma_k Tn_k + \delta_1 Tu_1 + \cdots + \delta_m Tu_m.\]
    Since $n_1, \ldots, n_k \in \nul T$, we have $Tn_i = 0$ for all $i = 1, \ldots, k$.
	Therefore,
    \[Tv = \delta_1 Tu_1 + \cdots + \delta_m Tu_m.\]
    Let $u = \delta_1 u_1 + \cdots + \delta_m u_m \in U$.
	Then $Tu = w$.
	Since $w$ was arbitrary in $\range T$, it follows that $\range T = \set{Tu \mid u \in U}$.
	Thus, we have found a subspace $U$ of $V$ such that $U \cap \nul T = \set{0}$ and $\range T = \set{Tu \mid u \in U}$.
\end{sol}

\begin{exercise}[3b12]
    Suppose $T$ is a linear map from $\f^4$ to $\f^2$ such that
    \[\nul T = \set{(x_1, x_2, x_3, x_4) \in \f^4 \mid x_1 = 5x_2 \text{ and } x_3 = 7x_4}.\]
    Prove that $T$ is surjective.
\end{exercise}

\begin{sol}
    We may write the null space of $T$ as
    \[\nul T = \set{(5x_2, x_2, 7x_4, x_4) \mid x_2, x_4 \in \f}.\]
    From this, we can see that $\nul T$ is spanned by the vectors $(5, 1, 0, 0)$ and $(0, 0, 7, 1)$.
	Therefore, $\dim \nul T = 2$.
	By the fundamental theorem of linear maps, we have
    \[\dim \f^4 = \dim \nul T + \dim \range T \implies 4 = 2 + \dim \range T.\]
    Solving for $\dim \range T$, we find that $\dim \range T = 2$.
	Since $\dim \f^2 = 2$, it follows that $\range T = \f^2$.
	Therefore, $T$ is surjective.
\end{sol}

\begin{exercise}[3b13]
    Suppose $U$ is a three-dimensional subspace of $\r^8$ and that $T$ is a linear map from $\r^8$ to $\r^5$ such that $\nul T = U$.
	Prove that $T$ is surjective.
\end{exercise}

\begin{sol}
    By the fundamental theorem of linear maps, we have
    \[\dim \r^8 = \dim \nul T + \dim \range T \implies 8 = 3 + \dim \range T.\]
    Solving for $\dim \range T$, we find that $\dim \range T = 5$.
	Since $\dim \r^5 = 5$, it follows that $\range T = \r^5$.
	Therefore, $T$ is surjective.
\end{sol}

\begin{exercise}[3b14]
    Prove that there does not exist a linear map from $\f^5$ to $\f^2$ whose null space equals $\set{(x_1, x_2, x_3, x_4, x_5) \in \f^5 \mid x_1 = 3x_2 \text{ and } x_3 = x_4 = x_5}$.
\end{exercise}

\begin{sol}
    We can express the null space as
    \[\nul T = \set{(3x_2, x_2, x_4, x_4, x_4) \mid x_2, x_4 \in \f}.\]
    From this, we can see that $\nul T$ is spanned by the vectors $(3, 1, 0, 0, 0)$ and $(0, 0, 1, 1, 1)$.
	Therefore, $\dim \nul T = 2$.
	By the fundamental theorem of linear maps, we have
    \[\dim \f^5 = \dim \nul T + \dim \range T \implies 5 = 2 + \dim \range T.\]
    Solving for $\dim \range T$, we find that $\dim \range T = 3$.
	However, $\dim \f^2 = 2$, which is less than $3$.
	This is a contradiction, as the dimension of the range cannot exceed the dimension of the target space.
	Therefore, no such linear map exists.
\end{sol}

\begin{exercise}[3b15]
    Suppose there exists a linear map on $V$ whose null space and range are both finite-dimensional.
	Prove that $V$ is finite-dimensional.

    \remark{The spanning list in the solution is actually a basis of $V$.}
\end{exercise}

\begin{sol}
    Let $T \in \lin{V}$ be such that $\nul T$ and $\range T$ are both finite-dimensional.
	Let $u_1, \ldots, u_m$ be a basis of $\nul T$, and let $w_1, \ldots, w_n$ be a basis of $\range T$.
	Select $v_j \in V$ be such that $Tv_j = w_j$ for $j = 1, \ldots, n$.
	We claim that the list $u_1, \ldots, u_m, v_1, \ldots, v_n$ spans $V$.
    
    Let $v \in V$.
	Then $Tv \in \range T$, so we may write
    \[Tv = \alpha_1 w_1 + \cdots + \alpha_n w_n\]
    for some scalars $\alpha_1, \ldots, \alpha_n \in \f$.
	Consider the vector
    \[u = v - (\alpha_1 v_1 + \cdots + \alpha_n v_n).\]
    Then
    \[Tu = Tv - (\alpha_1 Tv_1 + \cdots + \alpha_n Tv_n) = Tv - (\alpha_1 w_1 + \cdots + \alpha_n w_n) = 0.\]
    Therefore, $u \in \nul T$, and we can write
    \[u = \beta_1 u_1 + \cdots + \beta_m u_m\]
    for some scalars $\beta_1, \ldots, \beta_m \in \f$.
	Thus,
    \[v = u + (\alpha_1 v_1 + \cdots + \alpha_n v_n) = \beta_1 u_1 + \cdots + \beta_m u_m + \alpha_1 v_1 + \cdots + \alpha_n v_n,\]
    which shows that the list spans $V$.
	Hence, $V$ is finite-dimensional.
\end{sol}

\begin{exercise}[3b16]
    Suppose $V$ and $W$ are both finite-dimensional.
	Prove that there exists an injective linear map from $V$ to $W$ if and only if $\dim V \leq \dim W$.
\end{exercise}

\begin{sol}
    Suppose there exists an injective linear map $T \in \lin{V, W}$.
	By the fundamental theorem of linear maps, we have
    \[\dim V = \dim \nul T + \dim \range T.\]
    Since $T$ is injective, $\nul T = \set{0}$, so $\dim \nul T = 0$.
	Therefore,
    \[\dim V = 0 + \dim \range T = \dim \range T.\]
    Since $\range T$ is a subspace of $W$, we have $\dim \range T \leq \dim W$.
	Thus,
    \[\dim V = \dim \range T \leq \dim W.\]

    Conversely, suppose $\dim V \leq \dim W$.
	Let $n = \dim V$ and $m = \dim W$.
	Choose a basis $\set{v_1, \ldots, v_n}$ of $V$.
	Let $\set{w_1, \ldots, w_m}$ be a basis of $W$.
	Define a linear map $T \in \lin{V, W}$ by
    \[Tv_i = w_i \quad \text{for } i = 1, \ldots, n.\]
    Since $n \leq m$, we can assign each basis vector of $V$ to a distinct basis vector of $W$.
	To show that $T$ is injective, let $v = \alpha_1 v_1 + \cdots + \alpha_n v_n \in V$ such that $Tv = 0$.
	Then
    \[Tv = \alpha_1 Tv_1 + \cdots + \alpha_n Tv_n = \alpha_1 w_1 + \cdots + \alpha_n w_n = 0.\]
    Since $w_1, \ldots, w_n$ are linearly independent, it follows that $\alpha_1 = \cdots = \alpha_n = 0$.
	Therefore, $v = 0$, and $T$ is injective.
	Hence, there exists an injective linear map from $V$ to $W$ if and only if $\dim V \leq \dim W$.
\end{sol}

\begin{exercise}[3b17]
    Suppose $V$ and $W$ are both finite-dimensional.
	Prove that there exists a surjective linear map from $V$ onto $W$ if and only if $\dim V \geq \dim W$.
\end{exercise}

\begin{sol}
    Suppose there exists a surjective linear map $T \in \lin{V, W}$.
	By the fundamental theorem of linear maps, we have
    \[\dim V = \dim \nul T + \dim \range T.\]
    Since $T$ is surjective, $\range T = W$, so $\dim \range T = \dim W$.
	Therefore,
    \[\dim V = \dim \nul T + \dim W.\]
    Since $\dim \nul T \geq 0$, it follows that
    \[\dim V \geq \dim W.\]

    Conversely, suppose $\dim V \geq \dim W$.
	Let $n = \dim V$ and $m = \dim W$.
	Choose a basis $\set{v_1, \ldots, v_n}$ of $V$.
	Let $\set{w_1, \ldots, w_m}$ be a basis of $W$.
	Define a linear map $T \in \lin{V, W}$ by
    \[Tv_i = w_i \quad \text{for } i = 1, \ldots, m, \quad Tv_i = 0 \quad \text{for } i = m + 1, \ldots, n.\]
    To show that $T$ is surjective, let $w \in W$.
	Then $w$ can be expressed as a linear combination of the basis vectors of $W$:
    \[w = \beta_1 w_1 + \cdots + \beta_m w_m.\]
    Consider the vector $v = \beta_1 v_1 + \cdots + \beta_m v_m \in V$.
	Then
    \[Tv = \beta_1 Tv_1 + \cdots + \beta_m Tv_m = \beta_1 w_1 + \cdots + \beta_m w_m = w.\]
    Since $w$ was arbitrary in $W$, it follows that $\range T = W$.
	Therefore, $T$ is surjective.
	Hence, there exists a surjective linear map from $V$ onto $W$ if and only if $\dim V \geq \dim W$.
\end{sol}

\begin{exercise}[3b18]
    Suppose $V$ and $W$ are finite-dimensional and that $U$ is a subspace of $V$.
	Prove that there exists $T \in \lin{V, W}$ such that $\nul T = U$ if and only if $\dim U \geq \dim V - \dim W$.
\end{exercise}

\begin{sol}
    Suppose there exists $T \in \lin{V, W}$ such that $\nul T = U$.
	By the fundamental theorem of linear maps, we have
    \[\dim V = \dim \nul T + \dim \range T = \dim U + \dim \range T.\]
    Since $\range T$ is a subspace of $W$, we have $\dim \range T \leq \dim W$.
	Therefore,
    \[\dim V = \dim U + \dim \range T \leq \dim U + \dim W.\]
    Rearranging this inequality gives
    \[\dim U \geq \dim V - \dim W.\]

    Conversely, suppose $\dim U \geq \dim V - \dim W$.
	Let $n = \dim V$, $m = \dim W$, and $k = \dim U$.
	Let $u_1, \ldots, u_k$ be a basis of $U$.
	Extend this basis to a basis of $V$ by adding vectors $v_1, \ldots, v_{n - k}$.
	Define a linear map $T \in \lin{V, W}$ by
    \[Tu_i = 0 \quad \text{for } i = 1, \ldots, k, \quad Tv_j = w_j \quad \text{for } j = 1, \ldots, n - k,\]
    where $w_1, \ldots, w_{n - k}$ are linearly independent vectors in $W$.
	Since $\dim U = k$ and $\dim V - \dim W = n - m$, we have $k \geq n - m$, which implies that $n - k \leq m$.
	Therefore, we can choose $w_1, \ldots, w_{n - k}$ to be linearly independent in $W$.
	To show that $\nul T = U$, let $v \in \nul T$.
	Then $Tv = 0$.
	Express $v$ as a linear combination of the basis vectors of $V$:
    \[v = \alpha_1 u_1 + \cdots + \alpha_k u_k + \beta_1 v_1 + \cdots + \beta_{n - k} v_{n - k}.\]
    Applying $T$ to both sides, we have
    \[Tv = \alpha_1 Tu_1 + \cdots + \alpha_k Tu_k + \beta_1 Tv_1 + \cdots + \beta_{n - k} Tv_{n - k} = 0 + \beta_1 w_1 + \cdots + \beta_{n - k} w_{n - k} = 0.\]
    Since $w_1, \ldots, w_{n - k}$ are linearly independent, it follows that $\beta_1 = \cdots = \beta_{n - k} = 0$.
	Therefore,
    \[v = \alpha_1 u_1 + \cdots + \alpha_k u_k \in U.\]
    Hence, $\nul T \subseteq U$.
	Conversely, if $v \in U$, then $Tv = 0$ by the definition of $T$.
	Therefore, $U \subseteq \nul T$.
	Combining these two inclusions, we have $\nul T = U$.
	Thus, there exists $T \in \lin{V, W}$ such that $\nul T = U$ if and only if $\dim U \geq \dim V - \dim W$.
\end{sol}

\begin{exercise}[3b19]
    Suppose $W$ is finite-dimensional and $T \in \lin{V, W}$.
	Prove that $T$ is injective if and only if there exists $S \in \lin{W, V}$ such that $ST$ is the identity operator on $V$.
\end{exercise}

\begin{sol}
    Suppose $T$ is injective.
	Then $\nul T = \set{0}$.
	Also, $\range T$ is finite dimensional since it is a subspace of $W$.
	Since $T$ is injective, it is an isomorphism from $V$ onto $\range T$.
	Hence, $V$ is finite dimensional.
	Let $n = \dim V$ and $m = \dim W$.
	Since $T$ is injective, we have $\dim \range T = n$.
	Choose a basis $\set{v_1, \ldots, v_n}$ of $V$.
	Then $\set{Tv_1, \ldots, Tv_n}$ is linearly independent in $W$.
	Extend this set to a basis of $W$ by adding vectors $w_1, \ldots, w_{m - n}$.
	Define a linear map $S \in \lin{W, V}$ by
    \[S(Tv_i) = v_i \quad \text{for } i = 1, \ldots, n, \quad S(w_j) = 0 \quad \text{for } j = 1, \ldots, m - n.\]
    Then for any $v \in V$, we can express $v$ as a linear combination of the basis vectors:
    \[v = \alpha_1 v_1 + \cdots + \alpha_n v_n.\]
    Applying $T$ to both sides, we have
    \[Tv = \alpha_1 Tv_1 + \cdots + \alpha_n Tv_n.\]
    Now, applying $S$ to $Tv$, we get
    \[STv = S(\alpha_1 Tv_1 + \cdots + \alpha_n Tv_n) = \alpha_1 S(Tv_1) + \cdots + \alpha_n S(Tv_n) = \alpha_1 v_1 + \cdots + \alpha_n v_n = v.\]
    Therefore, $ST$ is the identity operator on $V$.

    Conversely, suppose there exists $S \in \lin{W, V}$ such that $ST$ is the identity operator on $V$.
	Let $v \in \nul T$.
	Then $Tv = 0$, and applying $S$ gives
    \[STv = S(0) = 0.\]
    Since $ST$ is the identity operator on $V$, we have $v = STv = 0$.
	Therefore, $\nul T = \set{0}$, which means $T$ is injective.
\end{sol}

\begin{exercise}[3b20]
    Suppose $W$ is finite-dimensional and $T \in \lin{V, W}$.
	Prove that $T$ is surjective if and only if there exists $S \in \lin{W, V}$ such that $TS$ is the identity operator on $W$.
\end{exercise}

\begin{sol}
    Suppose $T$ is surjective.
	Then $\range T = W$.
	Let $m = \dim W$ and choose a basis $\set{w_1, \ldots, w_m}$ of $W$.
	Since $T$ is surjective, for each $w_i$, there exists $v_i \in V$ such that $Tv_i = w_i$.
	Define a linear map $S \in \lin{W, V}$ by
    \[S(w_i) = v_i \quad \text{for } i = 1, \ldots, m.\]
    Then for any $w \in W$, we can express $w$ as a linear combination of the basis vectors:
    \[w = \alpha_1 w_1 + \cdots + \alpha_m w_m.\]
    Applying $S$ to both sides, we have
    \[Sw = S(\alpha_1 w_1 + \cdots + \alpha_m w_m) = \alpha_1 S(w_1) + \cdots + \alpha_m S(w_m) = \alpha_1 v_1 + \cdots + \alpha_m v_m.\]
    Now, applying $T$ to $Sw$, we get
    \[TSw = T(\alpha_1 v_1 + \cdots + \alpha_m v_m) = \alpha_1 Tv_1 + \cdots + \alpha_m Tv_m = \alpha_1 w_1 + \cdots + \alpha_m w_m = w.\]
    Therefore, $TS$ is the identity operator on $W$.

    Conversely, suppose there exists $S \in \lin{W, V}$ such that $TS$ is the identity operator on $W$.
	Let $w \in W$.
	Then
    \[TSw = w.\]
    This means that for every $w \in W$, there exists a vector $v = Sw \in V$ such that $Tv = T(Sw) = w$.
	Therefore, $\range T = W$, which means that $T$ is surjective.
\end{sol}

\begin{exercise}[3b21]
    Suppose $V$ is finite-dimensional, $T \in \lin{V, W}$, and $U$ is a subspace of $W$.
	Prove that $\set{v \in V \mid Tv \in U}$ is a subspace of $V$ and
    \[\dim\set{v \in V \mid Tv \in U} = \dim \nul T + \dim(U \cap \range T).\]
\end{exercise}

\begin{sol}
    Let $X = \set{v \in V \mid Tv \in U}$.
	We first show that $X$ is a subspace of $V$.
	Let $v_1, v_2 \in X$ and $\alpha, \beta \in \f$.
	Then
    \[T(\alpha v_1 + \beta v_2) = \alpha Tv_1 + \beta Tv_2.\]
    Since $Tv_1, Tv_2 \in U$ and $U$ is a subspace of $W$, it follows that $\alpha Tv_1 + \beta Tv_2 \in U$.
	Therefore, $\alpha v_1 + \beta v_2 \in X$, and thus $X$ is a subspace of $V$.

    Now, we will compute the dimension of $X$.
	Let $Y = U \cap \range T$.
	Let $\dim \null T = k$ and $\dim Y = m$.
	Let $n_1, \ldots, n_k$ be a basis of $\nul T$, and let $y_1, \ldots, y_m$ be a basis of $Y$.
	Since $y_i \in \range T$, there exist vectors $v_i \in V$ such that $Tv_i = y_i$ for $i = 1, \ldots, m$.
	We claim that the set $\set{n_1, \ldots, n_k, v_1, \ldots, v_m}$ is a basis of $X$.

    First, we show that this set spans $X$.
	Let $v \in X$.
	Then $Tv \in U$, and since $Tv \in \range T$, we have $Tv \in Y$.
	Therefore, there exist scalars $\alpha_1, \ldots, \alpha_m \in \f$ such that
    \[Tv = \alpha_1 y_1 + \cdots + \alpha_m y_m = \alpha_1 Tv_1 + \cdots + \alpha_m Tv_m.\]
    Let $v' = v - (\alpha_1 v_1 + \cdots + \alpha_m v_m)$.
	Then
    \[Tv' = Tv - T(\alpha_1 v_1 + \cdots + \alpha_m v_m) = Tv - (\alpha_1 Tv_1 + \cdots + \alpha_m Tv_m) = 0.\]
    Therefore, $v' \in \nul T$, and we can express $v'$ as a linear combination of the basis vectors of $\nul T$:
    \[v' = \beta_1 n_1 + \cdots + \beta_k n_k\]
    for some scalars $\beta_1, \ldots, \beta_k \in \f$.
	Thus,
    \[v = v' + (\alpha_1 v_1 + \cdots + \alpha_m v_m) = \beta_1 n_1 + \cdots + \beta_k n_k + \alpha_1 v_1 + \cdots + \alpha_m v_m.\]
    This shows that $v$ is a linear combination of the vectors in $\set{n_1, \ldots, n_k, v_1, \ldots, v_m}$, and hence this set spans $X$.

    Next, we show that this set is linearly independent.
	Suppose
    \[\gamma_1 n_1 + \cdots + \gamma_k n_k + \delta_1 v_1 + \cdots + \delta_m v_m = 0\]
    for some scalars $\gamma_1, \ldots, \gamma_k, \delta_1, \ldots, \delta_m \in \f$.
	Observe that $Tn_i = 0$ for all $i = 1, \ldots, k$, and $Tv_j = y_j$ for all $j = 1, \ldots, m$.
	Therefore,
    \[0 = T(\gamma_1 n_1 + \cdots + \gamma_k n_k + \delta_1 v_1 + \cdots + \delta_m v_m) = \delta_1 y_1 + \cdots + \delta_m y_m.\]
    Since $y_1, \ldots, y_m$ are linearly independent, it follows that $\delta_1 = \cdots = \delta_m = 0$.
	Therefore,
    \[\gamma_1 n_1 + \cdots + \gamma_k n_k = 0.\]
    Since $n_1, \ldots, n_k$ are linearly independent, $\gamma_1 = \cdots = \gamma_k = 0$.
	Thus, the set $\set{n_1, \ldots, n_k, v_1, \ldots, v_m}$ is linearly independent and a basis of $X$.
	Therefore,
    \[\dim X = k + m = \dim \nul T + \dim(U \cap \range T).\]
\end{sol}

\begin{exercise}[3b22]
    Suppose $U$ and $V$ are finite-dimensional vector spaces and $S \in \lin{V, W}$ and $T \in \lin{U, V}$.
	Prove that
    \[\dim \nul ST \leq \dim \nul S + \dim \nul T.\]
\end{exercise}

\begin{sol}
    Let $u \in U$ such that $u \in \nul ST$.
	Then $STu = 0$, which implies that $Tu \in \nul S$.
	Therefore, we can express $\nul ST$ as
    \[\nul ST = \set{u \in U \mid Tu \in \nul S}.\]
    Observe that this looks similar to the set in Exercise 3b21, where $T$ is a linear map from $U$ to $V$ and $\nul S$ is a subspace of $V$.
	By applying the result from Exercise 3b21, we have
    \[\dim \nul ST = \dim \nul T + \dim(\nul S \cap \range T).\]
    Since $\nul S \cap \range T$ is a subspace of $\nul S$, we have
    \[\dim(\nul S \cap \range T) \leq \dim \nul S.\]
    Therefore,
    \[\dim \nul ST = \dim \nul T + \dim(\nul S \cap \range T) \leq \dim \nul T + \dim \nul S. \qh\]
\end{sol}

\begin{exercise}[3b23]
    Suppose $U$ and $V$ are finite-dimensional vector spaces and $S \in \lin{V, W}$ and $T \in \lin{U, V}$.
	Prove that
    \[\dim \range ST \leq \min(\dim \range S, \dim \range T).\]
\end{exercise}

\begin{sol}
    Let $u \in U$ such that $STu \in \range ST$.
	Then $STu = S(Tu)$, which implies that $Tu \in \range T$.
	Therefore, we can express $\range ST$ as
    \[\range ST = \set{S(Tu) \mid u \in U}.\]
    Since $Tu \in \range T$, we can rewrite this as
    \[\range ST = \set{Sv \mid v \in \range T}.\]
    This shows that $\range ST$ is a subset of $\range S$, and hence
    \[\dim \range ST \leq \dim \range S.\]
    Now let $v_1, \ldots, v_m$ be a basis of $\range T$.
	Then for any $u \in U$, we can express $Tu$ as a linear combination of the basis vectors:
    \[Tu = \alpha_1 v_1 + \cdots + \alpha_m v_m\]
    for some scalars $\alpha_1, \ldots, \alpha_m \in \f$.
	Applying $S$ to both sides, we have
    \[STu = S(Tu) = S(\alpha_1 v_1 + \cdots + \alpha_m v_m) = \alpha_1 Sv_1 + \cdots + \alpha_m Sv_m.\]
    This shows that $\range ST$ is spanned by the vectors $Sv_1, \ldots, Sv_m$.
	Therefore,
    \[\dim \range ST \leq m = \dim \range T.\]
    Combining the two inequalities, we have
    \[\dim \range ST \leq \min(\dim \range S, \dim \range T). \qh\]
\end{sol}

\begin{exercise}[3b24]
    \begin{subexercises}
        \item Suppose $\dim V = 5$ and $S, T \in \lin{V}$ are such that $ST = 0$.
	Prove that $\dim \range TS \leq 2$.
        \item Give an example of $S, T \in \lin{\f^5}$ with $ST = 0$ and $\dim \range TS = 2$.
    \end{subexercises}
\end{exercise}

\begin{sole}
    \item Since $ST = 0$, then $S(Tv) = 0$ for all $v \in V$.
	This implies that $\range T \subseteq \nul S$.
	By the fundamental theorem of linear maps, we have
    \[5 = \dim \nul S + \dim \range S \geq \dim \range T + \dim \range S.\]
    Now, by \autoref{ex:3b23}, we have
    \[\dim \range TS \leq \min(\dim \range T, \dim \range S).\]
    Suppose that $\dim \range TS > 2$.
	Then both $\dim \range T$ and $\dim \range S$ must be greater than $2$, or equivalently, $\dim \range T \geq 3$ and $\dim \range S \geq 3$.
	This implies that
    \[\dim \range T + \dim \range S \geq 3 + 3 = 6 > 5,\]
    which is a contradiction.
	Therefore, $\dim \range TS \leq 2$.
    \item Let $v_1, v_2, v_3, v_4, v_5$ be a basis of $\f^5$.
	Define $S, T \in \lin{\f^5}$ by
    \[Sv_1 = v_1, \quad Sv_2 = v_2, \quad Sv_3 = Sv_4 = Sv_5 = 0,\]
    and
    \[Tv_1 = v_3, \quad Tv_2 = v_4, \quad Tv_3 = Tv_4 = Tv_5 = 0.\]
    Then $ST = 0$ since $S(Tv_i) = S(0) = 0$ for all $i = 1, \ldots, 5$.
	Now, we compute $\range TS$.
	We have
    \[TSv_1 = T(Sv_1) = Tv_1 = v_3, \quad TSv_2 = T(Sv_2) = Tv_2 = v_4, \quad TSv_3 = TSv_4 = TSv_5 = 0.\]
    Therefore, $\range TS = \spn(v_3, v_4)$, which has dimension $2$.
	Thus, we have constructed an example of $S, T \in \lin{\f^5}$ with $ST = 0$ and $\dim \range TS = 2$.
\end{sole}

\begin{exercise}[3b25]
    Suppose that $W$ is finite-dimensional and $S, T \in \lin{V, W}$.
	Prove that $\nul S \subseteq \nul T$ if and only if there exists $E \in \lin{W}$ such that $T = ES$.

    \remark{This exercise requires showing that some $R \in \lin{\range S, W}$ is well-defined.
	An issue arises when $T$ is not injective: for some $w \in \range S$, there may be distinct $v_1, v_2 \in V$ such that $Sv_1 = Sv_2 = w$.
	In this case, we must show that $Tv_1 = Tv_2$.}
\end{exercise}

\begin{sol}
    Suppose $\nul S \subseteq \nul T$.
	Define a linear map $R: \range S \to W$ by $R(Sv) = Tv$ for all $v \in V$.
	We need to show that $R$ is well-defined.
	Suppose $Sv_1 = Sv_2$ for some $v_1, v_2 \in V$.
	Then $S(v_1 - v_2) = 0$, which implies that $v_1 - v_2 \in \nul S$.
	Since $\nul S \subseteq \nul T$, we have $v_1 - v_2 \in \nul T$, which means that $T(v_1 - v_2) = 0$.
	Therefore, $Tv_1 = Tv_2$, and hence $R(Sv_1) = R(Sv_2)$.
	Thus, $R$ is well-defined.

    By \autoref{ex:3a13}, we can extend $R$ to some $E \in \lin{W}$.
	Then for any $v \in V$, we have
    \[Tv = R(Sv) = E(Sv),\]
    and hence $T = ES$.

    Conversely, suppose there exists $E \in \lin{W}$ such that $T = ES$.
	Let $v \in \nul S$.
	Then $Sv = 0$, and applying $E$ gives
    \[Tv = E(Sv) = E(0) = 0.\]
    Therefore, $v \in \nul T$, and hence $\nul S \subseteq \nul T$.
\end{sol}

\begin{exercise}[3b26]
    Suppose that $V$ is finite-dimensional and $S, T \in \lin{V, W}$.
	Prove that $\range S \subseteq \range T$ if and only if there exists $E \in \lin{V}$ such that $S = TE$.

    \remark{From the previous exercise, the reader may be tempted to show that a map is well-defined again.
	However, this is not necessary here, since we are defining a map on the entire space $V$, not just on $\range S$.}
\end{exercise}

\begin{sol}
    Suppose $\range S \subseteq \range T$.
	Let $v_1, \ldots, v_n$ be a basis of $V$.
	For each $i = 1, \ldots, n$, since $Sv_i \in \range S \subseteq \range T$, there exists $u_i \in V$ such that $Tu_i = Sv_i$.
	Define a linear map $E \in \lin{V}$ by
    \[Ev_i = u_i \quad \text{for } i = 1, \ldots, n.\]
    Then for any $v \in V$, we can express $v$ as a linear combination of the basis vectors:
    \[v = \alpha_1 v_1 + \cdots + \alpha_n v_n\]
    for some scalars $\alpha_1, \ldots, \alpha_n \in \f$.
	Applying $S$ and $TE$ to both sides, we have
    \[Sv = \alpha_1 Sv_1 + \cdots + \alpha_n Sv_n = \alpha_1 Tu_1 + \cdots + \alpha_n Tu_n = T(\alpha_1 u_1 + \cdots + \alpha_n u_n) = TEv.\]
    Therefore, $S = TE$.

    Conversely, suppose there exists $E \in \lin{V}$ such that $S = TE$.
	Let $w \in \range S$.
	Then there exists $v \in V$ such that $Sv = w$.
	Applying the definition of $S$, we have
    \[w = Sv = TE(v) = T(E(v)).\]
    Therefore, $w \in \range T$, and hence $\range S \subseteq \range T$.
\end{sol}

\begin{exercise}[3b27]
    Suppose $P \in \lin{V}$ and $P^2 = P$.
	Prove that $V = \nul P \op \range P$.
\end{exercise}

\begin{sol}
    Let $v \in V$.
	We can express $v$ as
    \[v = (v - Pv) + Pv.\]
    Observe that $Pv \in \range P$ by definition.
	Now, we need to show that $v - Pv \in \nul P$.
	Applying $P$ to $v - Pv$, we have
    \[P(v - Pv) = Pv - P^2v = Pv - Pv = 0.\]
    Therefore, $v - Pv \in \nul P$.
	This shows that every vector $v \in V$ can be expressed as a sum of a vector in $\nul P$ and a vector in $\range P$, which implies that
    \[V = \nul P + \range P.\]

    Next, we need to show that the sum is direct.
	Suppose $u \in \nul P \cap \range P$.
	Then there exists some $w \in V$ such that $u = Pw$.
	Since $u \in \nul P$, we have
    \[Pu = P(Pw) = P^2w = Pw = u.\]
    But since $u \in \nul P$, we also have $Pu = 0$.
	Therefore, we have
    \[u = Pu = 0.\]
    This shows that the intersection of $\nul P$ and $\range P$ is trivial, i.e., $\nul P \cap \range P = \set{0}$.
	Hence, the sum is direct, and we conclude that
    \[V = \nul P \oplus \range P. \qh\]
\end{sol}

\begin{exercise}[3b28]
    Suppose $D \in \lin{\pp(\r)}$ is such that $\deg Dp = (\deg p) - 1$ for every nonconstant polynomial $p \in \pp(\r)$.
	Prove that $D$ is surjective.
    
    \remark{Note that $\pp(\r)$ is infinite-dimensional, so we cannot use the fundamental theorem of linear maps directly.
	Instead, one should consider finite-dimensional subspaces of $\pp(\r)$.}
\end{exercise}

\begin{sol}
    Let $p \in \pp(\r)$ be an arbitrary polynomial.
	Consider the subspace $\pp_{n + 1}(\r)$ of polynomials of degree at most $n + 1$, where $n = \deg p$.
    Observe that no nonconstant polynomial can be in the null space of $D$, since nonconstant $q$ implies $\deg Dq = (\deg q) - 1 \geq 0$, hence $Dq \neq 0$.
	We now claim that $D1 = 0$.
	If it weren't, then let $\deg D1 = m$.
	Consider $D$ acting on $\pp_{m + 1}(\r)$.
	Since no nonconstant polynomial can be in $\nul D$, and $D1 \neq 0$ by assumption, we have $\nul D = \set 0$, which implies that $\dim \nul D = 0$.
	By the fundamental theorem of linear maps and the fact that $\dim \range D \leq \dim \pp_m(\r)$, we have
    \[\dim \nul D = \dim \pp_{m + 1}(\r) - \dim \range D \geq 1,\]
    which is a contradiction.
	Therefore, $D1 = 0$.
	By linearity, we have $Dc = cD1 = 0$ for all $c \in \r$.
	Hence, $\nul D = \spn\set 1$, which has dimension $1$.
	By the fundamental theorem of linear maps, we have
    \[\dim \range D = \dim \pp_{n + 1}(\r) - \dim \nul D = (n + 2) - 1 = n + 1 = \dim \pp_n(\r).\]
    Since $\dim \range D = \dim \pp_n(\r)$, it follows that $\range D = \pp_n(\r)$.
	Therefore, for any polynomial $p \in \pp_n(\r)$, there exists a polynomial $q \in \pp_{n + 1}(\r)$ such that $Dq = p$.
	Since $p$ was arbitrary, we conclude that $D$ is surjective.
\end{sol}

\begin{exercise}[3b29]
    Suppose $p \in \pp(\r)$.
	Prove that there exists a polynomial $q \in \pp(\r)$ such that $5q'' + 3q' = p$.

    \apnote{3b29}{ for an explicit method of determining $q$ given a system of equations.}
\end{exercise}

\begin{sol}
    Let $p \in \pp(\r)$ be an arbitrary polynomial.
	Define $D \in \lin{\pp_{n + 1}(\r), \pp_n(\r)}$ by $Dq = 5q'' + 3q'$, where $n = \deg p$.
	We want to show that $D$ is surjective.
	Observe that $D$ is well-defined since the degree of $5q'' + 3q'$ is at most $n$ for any polynomial $q \in \pp_{n + 1}(\r)$.

    Now suppose $q \in \nul D$ with $\deg q = m$.
	Then $\deg 3q' = m - 1$ and $\deg 5q'' = m - 2$.
	Therefore, $\deg Dq = \max(m - 1, m - 2) = m - 1$.
	Since $Dq = 0$, we must have $m - 1 < 0$, which implies that $m < 1$.
	Therefore, $q$ must be a constant polynomial.
	Hence, $\nul D$ consists of constant polynomials, which has dimension $1$.
	By the fundamental theorem of linear maps, we have
    \[\dim \range D = \dim \pp_{n + 1}(\r) - \dim \nul D = (n + 2) - 1 = n + 1 = \dim \pp_n(\r).\]
    Since $\dim \range D = \dim \pp_n(\r)$, it follows that $\range D = \pp_n(\r)$ so that $D$ is surjective.
	Since $p \in \range D$, there exists a polynomial $q \in \pp_{n + 1}(\r)$ such that $Dq = 5q'' + 3q' = p$.
\end{sol}

\begin{exercise}[3b30]
    Suppose $\varphi \in \lin{V, \f}$ and $\varphi \neq 0$.
	Suppose $u \in V$ is not in $\nul \varphi$.
	Prove that
    \[V = \nul \varphi \op \set{au \mid a \in \f}.\]
\end{exercise}

\begin{sol}
    Let $v \in V$.
	We want to show that $v$ can be expressed as a sum of a vector in $\nul \varphi$ and a scalar multiple of $u$.
	Since $\varphi(u) \neq 0$, we can define a scalar
    \[a = \frac{\varphi(v)}{\varphi(u)}.\]
    Then consider the vector
    \[w = v - au.\]
    We have
    \[\varphi(w) = \varphi(v - au) = \varphi(v) - a\varphi(u) = \varphi(v) - \frac{\varphi(v)}{\varphi(u)}\varphi(u) = 0.\]
    Therefore, $w \in \nul \varphi$.
	Thus, we can express $v$ as
    \[v = w + au,\]
    where $w \in \nul \varphi$ and $au$ is a scalar multiple of $u$.
	This shows that every vector in $V$ can be expressed as a sum of a vector in $\nul \varphi$ and a scalar multiple of $u$, which implies that
    \[V = \nul \varphi + \set{au \mid a \in \f}.\]

    Next, we need to show that the sum is direct.
	Suppose $w \in \nul \varphi \cap \set{au \mid a \in \f}$.
	Then there exists a scalar $a \in \f$ such that $w = au$.
	Since $w \in \nul \varphi$, we have
    \[0 = \varphi(w) = \varphi(au) = a\varphi(u).\]
    Since $\varphi(u) \neq 0$, it follows that $a = 0$.
	Therefore, $w = 0$, which shows that the intersection of $\nul \varphi$ and $\set{au \mid a \in \f}$ is trivial.
	Hence, the sum is direct, and we conclude that
    \[V = \nul \varphi \oplus \set{au \mid a \in \f}. \qh\]
\end{sol}

\begin{exercise}[3b31]
    Suppose $V$ is finite-dimensional, $X$ is a subspace of $V$, and $Y$ is a finite-dimensional subspace of $W$.
	Prove that there exists $T \in \lin{V, W}$ such that $\nul T = X$ and $\range T = Y$ if and only if $\dim X + \dim Y = \dim V$.
\end{exercise}

\begin{sol}
    Suppose there exists $T \in \lin{V, W}$ such that $\nul T = X$ and $\range T = Y$.
	By the fundamental theorem of linear maps, we have
    \[\dim V = \dim \nul T + \dim \range T = \dim X + \dim Y.\]

    Conversely, suppose $\dim X + \dim Y = \dim V$.
	Let $n = \dim V$, $k = \dim X$, and $m = \dim Y$.
	Then $n = k + m$.
	Choose a basis $\set{x_1, \ldots, x_k}$ of $X$ and extend it to a basis $\set{x_1, \ldots, x_k, v_1, \ldots, v_m}$ of $V$.
	Let $\set{y_1, \ldots, y_m}$ be a basis of $Y$.
	Define a linear map $T \in \lin{V, W}$ by
    \[Tx_i = 0 \quad \text{for } i = 1, \ldots, k,\]
    and
    \[Tv_j = y_j \quad \text{for } j = 1, \ldots, m.\]
    Then for any $v \in V$, we can express $v$ as a linear combination of the basis vectors:
    \[v = \alpha_1 x_1 + \cdots + \alpha_k x_k + \beta_1 v_1 + \cdots + \beta_m v_m\]
    for some scalars $\alpha_1, \ldots, \alpha_k, \beta_1, \ldots, \beta_m \in \f$.
	Applying $T$ to both sides gives
    \[Tv = \alpha_1 Tx_1 + \cdots + \alpha_k Tx_k + \beta_1 Tv_1 + \cdots + \beta_m Tv_m = 0 + \beta_1 y_1 + \cdots + \beta_m y_m.\]
    Therefore, $Tv \in Y$.
	Now suppose $y = \gamma_1 y_1 + \cdots + \gamma_m y_m \in Y$.
	Then
    \[T(\gamma_1 v_1 + \cdots + \gamma_m v_m) = \gamma_1 y_1 + \cdots + \gamma_m y_m = y.\]
    Hence, $\range T = Y$.
	Moreover, if $Tv = 0$, then $\beta_1 y_1 + \cdots + \beta_m y_m = 0$.
	Since $y_1, \ldots, y_m$ are linearly independent, it follows that $\beta_1 = \cdots = \beta_m = 0$.
	Thus, $v \in X$.
	Additionally, if $v \in X$, then $v = \alpha_1 x_1 + \cdots + \alpha_k x_k$ for some scalars $\alpha_1, \ldots, \alpha_k \in \f$, and hence $Tv = 0$.
	Therefore, $\nul T = X$.
	This shows that there exists a linear map $T \in \lin{V, W}$ such that $\nul T = X$ and $\range T = Y$.
\end{sol}

\begin{exercise}[3b32]
    Suppose $V$ is finite-dimensional with $\dim V > 1$.
	Show that if $\varphi: \lin{V} \to \f$ is a linear map such that $\varphi(ST) = \varphi(S)\varphi(T)$ for all $S, T \in \lin{V}$, then $\varphi = 0$.
    
    \note{The description of the two-sided ideals of $\lin{V}$ given by \autoref{ex:3a17} might be useful.}
\end{exercise}

\begin{sol}
    Observe that $\nul \varphi$ is a two-sided ideal of $\lin V$ as follows: if $X \in \nul \varphi$ and $Y \in \lin V$, then
    \[\varphi(XY) = \varphi(X)\varphi(Y) = 0 \cdot \varphi(Y) = 0, \quad \varphi(YX) = \varphi(Y)\varphi(X) = \varphi(Y) \cdot 0 = 0.\]
    Since $\dim V > 1$, then \autoref{ex:3a16} gives the existence of $S, T \in \lin V$ such that $ST \neq TS$.
	Then we have
    \[\varphi(ST - TS) = \varphi(S) \varphi(T) - \varphi(T) \varphi(S) = 0.\]
    Therefore, $ST - TS \in \nul \varphi$.
	Since $\nul \varphi$ is a two-sided ideal of $\lin V$, then by \autoref{ex:3a17}, we have $\nul \varphi = \lin V$.
	Hence, $\varphi = 0$.
\end{sol}

\begin{exercise}[3b33]
    Suppose that $V$ and $W$ are real vector spaces and $T \in \lin{V, W}$.
	Define $T_\c: V_\c \to W_\c$ by
    \[T_\c(u + iv) = Tu + iTv\]
    for all $u, v \in V$.
    \begin{subexercises}
        \item Show that $T_\c$ is a (complex) linear map from $V_\c$ to $W_\c$.
        \item Show that $T_\c$ is injective if and only if $T$ is injective.
        \item Show that $\range T_\c = W_\c$ if and only if $\range T = W$.
    \end{subexercises}
    
    \note{See Exercise 8 in Section 1B for the definition of the complexification $V_\c$.
	The linear map $T_\c$ is called the \textbf{complexification} of the linear map $T$.}
\end{exercise}

\begin{sole}
    \item Let $u + vi$ and $w + xi \in V_\c$.
	Then for any $a + bi \in \c$, we have
    \begin{align*}
        T_\c((a + bi)(u + vi) + (w + xi)) & = T_\c((au - bv + (av + bu)i) + (w + xi)) \\
        & = T(au - bv + w) + iT(av + bu + x) \\
        & = aTu - bTv + Tw + i(aTv + bTu + Tx) \\
        & = (a + bi)(Tu + iTv) + (Tw + iTx) \\
        & = (a + bi)T_\c(u + vi) + T_\c(w + xi).
    \end{align*}
    Therefore, $T_\c$ is a complex linear map from $V_\c$ to $W_\c$.
    \item Suppose $T_\c$ is injective.
	Let $v \in V$ be such that $Tv = 0$.
	Then
    \[T_\c(v + i0) = Tv + iT0 = 0 + i0 = 0.\]
    Since $T_\c$ is injective, we have $v + i0 = 0$, which implies that $v = 0$.
	Therefore, $T$ is injective.

    Conversely, suppose $T$ is injective.
	Let $u + iv \in V_\c$ be such that $T_\c(u + iv) = 0$.
	Then
    \[Tu + iTv = 0.\]
    This implies that $Tu = 0$ and $Tv = 0$.
	Since $T$ is injective, we have $u = 0$ and $v = 0$.
	Therefore, $u + iv = 0$, and hence $T_\c$ is injective.
    \item Suppose $\range T_\c = W_\c$.
	Let $w \in W$.
	Then $w + i0 \in W_\c$, and since $\range T_\c = W_\c$, there exists $u + iv \in V_\c$ such that
    \[T_\c(u + iv) = Tu + iTv = w + i0.\]
    This implies that $Tu = w$ and $Tv = 0$.
	Therefore, $w \in \range T$, and hence $\range T = W$.

    Conversely, suppose $\range T = W$.
	Let $w_1 + iw_2 \in W_\c$.
	Since $\range T = W$, there exist $u, v \in V$ such that $Tu = w_1$ and $Tv = w_2$.
	Then
    \[T_\c(u + iv) = Tu + iTv = w_1 + iw_2.\]
    Therefore, $w_1 + iw_2 \in \range T_\c$, and hence $\range T_\c = W_\c$.
\end{sole}